\documentclass{amsart}
% For dealing with references we use the comment environment
\usepackage{verbatim}
\newenvironment{reference}{\comment}{\endcomment}
%\newenvironment{reference}{}{}
\newenvironment{slogan}{\comment}{\endcomment}
\newenvironment{history}{\comment}{\endcomment}

% For commutative diagrams we use Xy-pic
\usepackage[all]{xy}

% We use 2cell for 2-commutative diagrams.
\xyoption{2cell}
\UseAllTwocells

% We use multicol for the list of chapters between chapters
\usepackage{multicol}

% This is generally recommended for better output
\usepackage{lmodern}
\usepackage[T1]{fontenc}

% For cross-file-references

% Package for hypertext links:
\usepackage{hyperref}

% For any local file, say "hello.tex" you want to link to please

% Theorem environments.
%
\theoremstyle{plain}
\newtheorem{theorem}[subsection]{Theorem}
\newtheorem{proposition}[subsection]{Proposition}
\newtheorem{lemma}[subsection]{Lemma}

\theoremstyle{definition}
\newtheorem{definition}[subsection]{Definition}
\newtheorem{example}[subsection]{Example}
\newtheorem{exercise}[subsection]{Exercise}
\newtheorem{situation}[subsection]{Situation}

\theoremstyle{remark}
\newtheorem{remark}[subsection]{Remark}
\newtheorem{remarks}[subsection]{Remarks}

\numberwithin{equation}{subsection}

% Macros
%
\def\lim{\mathop{\mathrm{lim}}\nolimits}
\def\colim{\mathop{\mathrm{colim}}\nolimits}
\def\Spec{\mathop{\mathrm{Spec}}}
\def\Hom{\mathop{\mathrm{Hom}}\nolimits}
\def\Ext{\mathop{\mathrm{Ext}}\nolimits}
\def\SheafHom{\mathop{\mathcal{H}\!\mathit{om}}\nolimits}
\def\SheafExt{\mathop{\mathcal{E}\!\mathit{xt}}\nolimits}
\def\Sch{\mathit{Sch}}
\def\Mor{\mathop{\mathrm{Mor}}\nolimits}
\def\Ob{\mathop{\mathrm{Ob}}\nolimits}
\def\Sh{\mathop{\mathit{Sh}}\nolimits}
\def\NL{\mathop{N\!L}\nolimits}
\def\CH{\mathop{\mathrm{CH}}\nolimits}
\def\proetale{{pro\text{-}\acute{e}tale}}
\def\etale{{\acute{e}tale}}
\def\QCoh{\mathit{QCoh}}
\def\Ker{\mathop{\mathrm{Ker}}}
\def\Im{\mathop{\mathrm{Im}}}
\def\Coker{\mathop{\mathrm{Coker}}}
\def\Coim{\mathop{\mathrm{Coim}}}

% Boxtimes
%
\DeclareMathSymbol{\boxtimes}{\mathbin}{AMSa}{"02}

%
% Macros for moduli stacks/spaces
%
\def\QCohstack{\mathcal{QC}\!\mathit{oh}}
\def\Cohstack{\mathcal{C}\!\mathit{oh}}
\def\Spacesstack{\mathcal{S}\!\mathit{paces}}
\def\Quotfunctor{\mathrm{Quot}}
\def\Hilbfunctor{\mathrm{Hilb}}
\def\Curvesstack{\mathcal{C}\!\mathit{urves}}
\def\Polarizedstack{\mathcal{P}\!\mathit{olarized}}
\def\Complexesstack{\mathcal{C}\!\mathit{omplexes}}
% \Pic is the operator that assigns to X its picard group, usage \Pic(X)
% \Picardstack_{X/B} denotes the Picard stack of X over B
% \Picardfunctor_{X/B} denotes the Picard functor of X over B
\def\Pic{\mathop{\mathrm{Pic}}\nolimits}
\def\Picardstack{\mathcal{P}\!\mathit{ic}}
\def\Picardfunctor{\mathrm{Pic}}
\def\Deformationcategory{\mathcal{D}\!\mathit{ef}}

\setlength{\emergencystretch}{3em}
\begin{document}
\title[Stacks Project, Tag 0F32]{494 --- Locally quasi-finite componentwise dominant morphisms to geometrically unibranch locally Noetherian schemes are universally open}
\maketitle
\noindent Copyright (C) 2005--2025 Johan de Jong. Stacks Project authors. Source: \href{https://stacks.math.columbia.edu/tag/0F32}{Tag 0F32}.

\noindent Editorial difficulty note (not author text). Difficulty H1: The proof reduces universal openness to lifting generalizations and isolates a finite component after elementary étale base change. Geometric unibranchness forces the component image to fill the local target, so closedness of the finite map and uniqueness of its special point produce the required lift..

\noindent Editorial provenance note (not author text). History: excerpt from revision \texttt{a0444\allowbreak 6e57e\allowbreak c1fbc\allowbreak 252a8\allowbreak 71afc\allowbreak ec775\allowbreak 2fb28\allowbreak 07b14}, more-morphisms.tex, lines 21787--21844. Reference presentation was adapted; original-author context and core support blocks were retained or added. The mathematical wording is unchanged. Licensed under GNU FDL 1.2 or later; no invariant sections or cover texts. See ../licenses/COPYING. Context blocks reproduce author definitions and supporting results; they are not additional counted problems.

\begin{lemma}
\label{lemma-quasi-finite-Noetherian-universally-open}
Let $f : X \to Y$ be a morphism of schemes. If
\begin{enumerate}
\item $f$ is locally quasi-finite,
\item $Y$ is geometrically unibranch and locally Noetherian, and
\item every irreducible component of $X$ dominates
an irreducible component of $Y$,
\end{enumerate}
then $f$ is universally open.
\end{lemma}

\begin{proof}
For any $n$ the scheme $\mathbf{A}^n \times Y$ is geometrically unibranch
by Lemma \href{https://stacks.math.columbia.edu/tag/0DQ2}{lemma number of branches and smooth (Tag 0DQ2)} and
Properties, Lemma \href{https://stacks.math.columbia.edu/tag/0C39}{lemma number of branches 1 (Tag 0C39)}. Hence
the hypotheses of the lemma hold for the morphisms
$\mathbf{A}^n \times X \to \mathbf{A}^n \times Y$ for all $n$.
By Lemma \href{https://stacks.math.columbia.edu/tag/0F31}{lemma test universally open (Tag 0F31)}
it suffices to prove $f$ is open. By Morphisms, Lemma
\href{https://stacks.math.columbia.edu/tag/01U1}{lemma locally finite presentation universally open (Tag 01U1)}
it suffices to show that generalizations lift along $f$.
Suppose that $y' \leadsto y$ is a specialization of points in
$Y$ and $x \in X$ is a point mapping to $y$.
As in Lemma \href{https://stacks.math.columbia.edu/tag/02LK}{lemma etale makes quasi finite finite at point (Tag 02LK)}
choose a diagram
$$
\xymatrix{
u \ar[d] & U \ar[d] \ar[r] & X \ar[d] \\
v & V \ar[r] & Y
}
$$
where $(V, v) \to (Y, y)$ is an elementary \'etale neighbourhood,
$U \to V$ is finite, $u$ is the unique point of $U$ mapping to $v$,
$U \subset V \times_Y X$ is open, and $v \mapsto y$ and $u \mapsto x$.
Let $E$ be an irreducible component of $U$ passing through $u$
(there is at least one of these). Since $U \to X$ is \'etale, $E$
maps to an irreducible component of $X$,
which in turn dominates an irreducible component of $Y$ (by assumption).
Since $U \to V$ is finite hence closed, we conclude that
the image $E' \subset V$ of $E$ is an irreducible closed subset
passing through $v$ which dominates an irreducible component of $Y$.
Since $V \to Y$ is \'etale $E'$ must be an irreducible component
of $V$ passing through $v$.
Since $Y$ is geometrically unibranch we see that $E'$ is the unique irreducible
component of $V$ passing through $v$ (Lemma \href{https://stacks.math.columbia.edu/tag/0CB4}{lemma nr branches (Tag 0CB4)}).
Since $V$ is locally Noetherian we may after shrinking $V$
assume that $E' = V$ (equality of sets).

\medskip\noindent
Since $V \to Y$ is \'etale we can find a specialization
$v' \leadsto v$ whose image is $y' \leadsto y$.
By the above we can find $u' \in U$ mapping to $v'$.
Then $u' \leadsto u$ because $u$ is the only point of
$U$ mapping to $v$ and $U \to V$ is closed.
Then finally the image $x' \in X$ of $u'$ is a point
specializing to $x$ and mapping to $y'$ and the proof is complete.
\end{proof}
\end{document}
